\section{Limitations}

Despite the positive empirical findings, several limitations of the current study must be acknowledged:

First, graph-based recommendation architectures introduce significant computational and memory overhead during training and inference. The dual-graph convolution layers in CombiGCN require maintaining and propagating messages over large item-item similarity matrices. Similarly, end-to-end representation learning and modal projection in BM3 and FREEDOM demand high GPU memory footprints. As a result, scaling these GNN models to real-world corporate catalogs—which routinely contain tens of millions of active users, outfits, and transactions—is constrained by memory boundaries and training time limitations.

Second, the reliance on offline feature extraction represents a static bottleneck. Visual and textual features are pre-extracted using pre-trained networks (MobileNetV2, CLIP, and BERT) and cached as static embeddings. This design prevents the system from capturing real-time temporal dynamics, such as rapid shifts in fashion trends, changing user style preferences, or seasonal transition cycles (e.g., winter coat versus summer dress recommendations).

Third, the empirical conclusions are subject to evaluation-protocol constraints. Results are reported from a single training run per configuration, without multiple random seeds or statistical significance testing. Given the small test set (an average of $3.81$ ground-truth items per user) and the near-floor absolute metric values, the narrow margins between the top configurations (e.g., BM3 versus CombiGCN) should be interpreted as indicative trends rather than statistically confirmed differences. Furthermore, this study did not include a non-personalized popularity baseline, did not separately benchmark a standalone LightGCN interaction-only model, and fixed the embedding dimension at $d=512$ without tuning; consequently, the magnitude of the gains attributable to multimodal fusion and the sensitivity of the models to embedding capacity remain to be verified. These protocol extensions are left for future work.

Finally, while multimodal graphs alleviate data sparsity, extreme interaction sparsity boundaries remain a challenge. For extremely casual users or tail items with fewer than three initial interactions, the collaborative and semantic signals propagate poorly along graph neighborhoods, which limits the model's capacity to deliver highly personalized recommendations under extreme sparsity constraints.
